\documentclass[10pt]{article}
\usepackage[margin=20mm]{geometry}
\usepackage{amsmath,amssymb,mathtools,fontspec,unicode-math,xurl,hyperref}
\IfFontExistsTF{DejaVu Serif}{\setmainfont{DejaVu Serif}}{\setmainfont{Latin Modern Roman}}
\setmathfont{latinmodern-math.otf}
\setlength{\parindent}{0pt}\setlength{\parskip}{6pt}\sloppy
\title{Technical Verification of
Singularities of Ricci Flow and Diffeomorphisms}
\author{AI-assisted technical review}
\date{}
\begin{document}
\maketitle
Abstract: The acquired publisher version is checked through the weighted elliptic identities, frequency estimates, propagation of almost splitting, quadratic rigidity and gauge iteration. A core propagation argument invokes a relative endpoint estimate without establishing its required hypothesis. A scalar diagnostic confirms that the displayed input bounds alone cannot give the asserted decay. Retaining a boundary term yields a weaker estimate on a smaller expanded region; its positive exponential rate is sufficient for the subsequent qualitative iteration. Separately, a Gaussian Hermite polynomial disproves the precise integrated estimate (4.5), whose correction preserves the polynomial degree. These are respectively a Major proof defect and a Minor quantitative defect. Neither is a counterexample to cylindrical shrinker rigidity. Scientific adjudication remains pending.\par
\section*{Paper identity and inspected scope}
The paper is Tobias Holck Colding and William P. Minicozzi II, Singularities of Ricci flow and diffeomorphisms, Publications Mathématiques de l’IHÉS 142 (2025), 75–152, DOI 10.1007/s10240-025-00157-1 (https://doi.org/10.1007/s10240-025-00157-1) [paper]. The reviewed source is the acquired 78-page publisher version of record, published online on 22 September 2025. Physical PDF page \(p\) corresponds to printed page \(p+74\). The separately acquired arXiv:2109.06240v4 author source, dated 19 September 2025, is a distinct version and is not asserted identical to the publisher PDF.\par
The bounded checks cover Sections 2–5, Sections 6, 8–9, and Sections 7, 10 and the reference list, with overlapping root checks of the operative statements and proof interfaces. Critical source pages were actually viewed, separately from text reading and rendering. This is an initial technical review, not independent acceptance or a complete reproof of every cited foundation. No separate supplement or original code repository was identified in the bounded availability checks; absence of code is not a mathematical defect.\par
\section*{Definitions, quantifiers and proof architecture}
The normalized shrinker equation and weight identities are \[\operatorname{Ric}+\operatorname{Hess}f=\tfrac12g,\qquad
 S+|\nabla f|^2=f,\qquad \Delta f+S=\tfrac n2.\] The norms below are weighted by \(e^{-f}\), unless otherwise specified. The drift operator is \(\mathcal L=\Delta-\nabla_{\nabla f}\), while \(\mathcal P=\operatorname{div}_f\operatorname{div}_f^*\) is nonnegative. For the growth estimates, \(f\) is proper, \(S\ge0\), and \(b=2\sqrt f\). For a tensor \(u\), the relevant quantities are \[\begin{aligned}
 I(r)&=r^{1-n}\int_{b=r}|u|^2|\nabla b|,\\
 D(r)&=\frac{r^{2-n}}2\int_{b=r}
       \left\langle\nabla|u|^2,\frac{\nabla b}{|\nabla b|}\right\rangle,
 \qquad U(r)=D(r)/I(r).\end{aligned}\] The solid-integral definition of \(D\) agrees with this boundary formula at regular radii. Local absolute continuity is used when integrating the identities.\par
The model is \(\Sigma=N^\ell\times\mathbb R^{n-\ell}\), with \(N\) closed, \(\operatorname{Ric}_N=g_N/2\), and \(\bar f=|x|^2/4+\ell/2\). Condition \((\star)\) in Section 6 is a spectral gap on trace-free \(N\)-tensors with a constant \(C_{N,n}\). It is not a statement that every positive Einstein manifold is admissible or that all constants are uniform over all \(N\). For round shrinking spheres the operative normalization is radius \(\sqrt{2(\ell-1)}\), \(\ell\ge2\), as fixed by the Ricci equation and curvature formula.\par
Theorem 7.1 extends the strong weighted closeness \((\star_R)\) to the weaker uniform closeness \((\dagger_{(1+\beta)R})\). Theorem 10.1 improves \((\dagger_\rho)\) to \((\star_{\theta\rho})\), for any fixed \(0<\theta<1\). Choosing \(\theta(1+\beta)>1\) gives expanding scales. The precise tangent-flow result is Theorem 10.8, which assumes the compactness and shrinking-limit conclusions of Mantegazza–Müller Theorem 1.4 [mm]; closed Type I flows are an included case. The unrestricted introductory wording is not substituted for those operative hypotheses. In the exact author version consulted, Theorem 1.4(2) allows arbitrary blow-up sequences with a fixed adjoint heat-kernel potential. It therefore supplies subsequences at the exit times used in Section 10; continuity of a time-dependent choice of entropy minimizer is not required. The potential normalization must also be aligned by its additive entropy constant. The common limiting entropy discussed there permits the same shift for the compared limits. The noncompact extension mentioned in its remarks is not treated here as a proved unrestricted theorem.\par
\section*{Substantive findings}
Problem 1 (The relative endpoint hypothesis is missing in the core propagation step).  The deduction of (7.20) from Proposition 4.16, in the proof of Proposition 7.2, is invalid from the bounds established there. This is a Major defect in the displayed derivation of almost splitting.\par
\subsection*{Original claim and exact location}
Equation reference: eq:relative. Proposition 4.16, PDF pp. 34–35 (printed pp. 108–109), assumes \(\langle\mathcal Lu,u\rangle\ge-\lambda|u|^2\) on \(r_1\le b\le L\), with \(r_1\) sufficiently large, and the relative endpoint condition \[
 D(L)\le \exp\!\left(\frac{2L-1}{6}\right)I(r_1).\] It then bounds \(I(r)\) by a fixed power of \(r/r_1\) times \(I(r_1)\) for \(r\le L-1\).\par
In the proof of Proposition 7.2, PDF p. 49 (printed p. 123), set \(u=\operatorname{Hess}v_i\) and \(L=(1+\alpha_0)R-R^{-1}\). The displayed inputs are a polynomial bound on \(I_u,D_u\) up to \(L\), the local differential inequality (7.19), and \[I_u(s)\le C_s e^{-R^2/8}
 \quad\hbox{for each fixed sufficiently large }s.\] The text then invokes Proposition 4.16 to assert \[I_u(r)\le CR^m e^{-R^2/8}\qquad(r\le L-1),\] and derives the pointwise Hessian estimate (7.20). A polynomial upper bound on \(D_u(L)\) and a small upper bound on \(I_u(s)\) do not establish [eq:relative] (\#eq:relative). In particular, no lower bound on \(I_u(s)\), or corresponding small relative bound on \(D_u(L)\), is supplied.\par
\subsection*{An exact diagnostic of the insufficient inputs}
This diagnostic addresses that inference, not all geometric hypotheses of Proposition 7.2. On the one-dimensional Gaussian shrinker, let \[f(x)=x^2/4,\qquad
 w_L(x)=\chi_L(x)e^{(x^2-L^2)/4},\] where \(0\le\chi_L\le1\) is smooth, equals one on \([-L,L]\), and vanishes outside \([-L-1,L+1]\). On every interval \(s\le |x|\le L\), \[\mathcal Lw_L=\tfrac12w_L,\qquad
 I(r)=2e^{(r^2-L^2)/2},\qquad
 D(r)=r^2e^{(r^2-L^2)/2}.\] Thus the required differential inequality holds for any positive \(\lambda\), and \(I,D\) are polynomially bounded up to \(L\). At any fixed inner radius \(s\) exceeding the required threshold, \(I(s)\) is exponentially small in \(L^2\), but \[D(L)=L^2,\qquad I(L-1)=2e^{-L+1/2}.\] The ratio \(D(L)/I(s)\) grows as \(e^{L^2/2}\), so [eq:relative] (\#eq:relative) fails for all sufficiently large \(L\).\par
Even the small global weighted norm does not resolve this inference. With \(L=11R/10\), \[\begin{aligned}
 \|w_L\|_{L^2(e^{-f})}^2
 &\le 2(L+1)e^{-L^2/4+L/2+1/4}\\
 &\le e^{-L^2/4+3L/2+5/4}
 \le e^{-R^2/8}\quad(R\ge20).\end{aligned}\] Nevertheless \(I(L-1)\) cannot be bounded by \(CR^m e^{-R^2/8}\) for any fixed \(C,m\) as \(R\to\infty\). This is a counterexample to the proposed implication between the listed analytic bounds. The cutoff tensor is not asserted to be the Hessian of a global low eigenfunction on an admissible perturbed shrinker. It therefore does not disprove Proposition 7.2 itself or the main rigidity theorem.\par
\subsection*{A proved boundary-term correction}
Equation reference: eq:boundary. The proof of Proposition 4.16 supplies a useful replacement. Fix \(0<\delta<\min\{1,2\lambda\}\), set \[a=2\lambda+\delta,\qquad q=2a,\qquad K=D-aI,\] and choose \(s\) above the fixed threshold in the differential estimates. The same inequality used in that proof gives \(K'(t)\ge tK(t)/3\) wherever \(K(t)>0\). Absolute continuity implies that a positive value cannot cross down through zero. Consequently \[K_+(t)\le e^{-(L^2-t^2)/6}D_+(L),\qquad s\le t\le L.\] Combining this with \(I'\le2D/t=qI/t+2K/t\) and integrating gives \[\begin{aligned}
 I(r)&\le \left(\frac rs\right)^q I(s)
 +2r^q D_+(L)\int_s^r
       t^{-q-1}e^{-(L^2-t^2)/6}\,dt\\
 &\le \left(\frac rs\right)^q
 \left\{I(s)+\frac2qD_+(L)e^{-(L^2-r^2)/6}\right\}.
 \end{aligned}\] This proof does not assume [eq:relative] (\#eq:relative). It exposes the missing boundary contribution.\par
At \(r=L-1\), that term has only \(e^{-O(L)}\) decay. Instead restrict to \(r\le(1+\alpha_0/2)R\). For large \(R\), the established polynomial endpoint bound and [eq:boundary] (\#eq:boundary) give \[I(r)\le CR^M e^{-cR^2},\qquad
 0<c\le\min\!\left\{\frac18,
 \frac{(1+\alpha_0)^2-(1+\alpha_0/2)^2}{12}\right\}.\] Thus a smaller expansion factor and a model-dependent positive exponential rate can be recovered. This does not recover the printed fixed rates \(1/8\) and \(1/16\) on their stated domains.\par
\subsection*{Impact, counterevidence and repair range}
Proposition 7.2 supplies the almost parallel fields used throughout Theorem 7.1, which is one of the two iteration steps in Theorem 10.2. The missing hypothesis is therefore tied directly to the core contribution. Its severity concerns that demonstrated derivation, not a claim that the qualitative theorem is false.\par
There is substantial positive counterevidence to a broader failure. Local elliptic estimates, square roots, interpolation and integration of the almost parallel fields lose only polynomial factors and fixed fractions of a positive exponent. On a further reduced expansion factor \(\beta>0\), they give \(CR^M e^{-c_1R^2}\) uniform errors for some \(c_1>0\). Choosing \(a_0<c_1/[2(1+\beta)^2]\) and enlarging the starting radius gives the \((\dagger_{(1+\beta)R})\) needed by Theorem 7.1. Section 10 requires only \(a_0>0\), with no positive lower threshold.\par
For completeness, the apparent global flow/completeness assertion on p. 51 can be localized. Let \(G_{ij}=\langle\nabla u_i,\nabla u_j\rangle\) and \(X_i=\sum_j(G^{-1})_{ij}\nabla u_j\), so \(du_j(X_i)=\delta_{ij}\). Flowing from the central zero section along \(\sum_i x_iX_i\) gives \(u(H(q,tx))=tx\) exactly. The almost parallel estimates imply \(f(H(q,tx))=f(q)+t^2|x|^2/4+O(\varepsilon R^2)\). Choose the target radius strictly below the coordinate boundary and that boundary strictly inside the controlled region. The flow remains inside that region, the restricted map is a proper local diffeomorphism, and its degree is one because its zero-section fiber is unique. This supplies the required local product map without assuming the uncontrolled exterior already has product topology.\par
Finally, in the improvement step choose \[\frac1{1+\beta}<\theta<1,\qquad
 0<\epsilon<\min\{1/2,\,2(1-\theta^2)\}.\] Then \(1-\epsilon/2>\theta^2\). The Gaussian shell errors in Theorem 9.1 dominate all polynomial losses and give \((\star_{\theta\rho})\); the new scale multiplier is \(\theta(1+\beta)>1\). The weaker positive-rate repair therefore supports the qualitative extension/improvement mechanism. The original sharp propagation rate remains unverified, and no author adoption of this repair is known.\par
Problem 2 (The integrated frequency bound has the wrong exponential coefficient).  Theorem 4.1, equation (4.5), PDF p. 21 (printed p. 95), is false as printed. This is a Minor quantitative defect.\par
\subsection*{Claim, hypotheses and complete counterexample}
Equation reference: eq:printed. The theorem assumes \(u,\mathcal Lu\in L^2(e^{-f})\), the weight identities, \(\langle\mathcal Lu,u\rangle\ge-\lambda|u|^2\), and that \(u\) does not vanish identically outside a compact set. It states, for every \(\epsilon>0\) and all sufficiently large \(r_2>r_1\), \[
 I(r_2)\le I(r_1)(r_2/r_1)^{4\lambda}
 e^{(\bar\mu+\epsilon)(r_1^{-2}-r_2^{-2})},
 \qquad \bar\mu=2n+4\max\{\lambda-1,0\}.\]\par
Take \(M=\mathbb R^n\), \(f=|x|^2/4\), \(S=0\), and \(u=|x|^2-2n\). All structural identities hold, \(f\) is smooth and proper, and \(\mathcal Lu=-u\). Both \(u\) and \(\mathcal Lu\) are Gaussian square integrable; \(u\) is nonzero outside a compact set. Thus \(\lambda=1\) and \[I(r)=|\mathbb S^{n-1}|(r^2-2n)^2.\] Choose \(0<\epsilon<2n\), \(r_2=2r_1=2r\), and \(r^2>2n\). With \(a=2n/r^2\in(0,1)\), \[\begin{aligned}
 \log\frac{I(2r)}{16I(r)}
 &=2\int_{a/4}^a\frac{dt}{1-t}
 >\frac{3n}{r^2},\\
 (\bar\mu+\epsilon)(r^{-2}-(2r)^{-2})
 &=\frac{3(2n+\epsilon)}{4r^2}
 <\frac{3n}{r^2}.\end{aligned}\] Hence [eq:printed] (\#eq:printed) fails for arbitrarily large \(r\); no choice of the theorem’s threshold removes the example. The exact rational diagnostic file additionally certifies 24 finite instances using \(e^A<1/(1-A)\), \(0<A<1\), rather than floating-point comparison.\par
\subsection*{Correction and bounded impact}
Lemma 4.7 gives \((\log I)'\le2U/r\). Integrating the preceding correct bound \(U(r)\le2\lambda[1+(\bar\mu+\epsilon)/r^2]\) yields \[I(r_2)\le I(r_1)(r_2/r_1)^{4\lambda}
 e^{\,2\lambda(\bar\mu+\epsilon)(r_1^{-2}-r_2^{-2})}.\] The missing factor is \(2\lambda\). The original image and the proof on p. 30 confirm that this is not an extraction error. For \(\lambda=1/2\) the coefficients coincide, and for other parameters the correction does not change the polynomial degree \(4\lambda\). Later estimates with a fixed multiplicative constant, or a positive power loss, can absorb it by increasing the starting radius. This defect does not itself damage the qualitative propagation or rigidity argument.\par
\section*{Other checked mechanisms and honest limits}
The principal symbol of \(\mathcal P\) is \(\tfrac12(|\xi|^2\operatorname{Id}+\xi\otimes\xi)\), hence positive definite. Weighted cutoff arguments control the boundary terms at infinity; properness and the weighted concentration estimate combine with local Rellich compactness to give the discrete spectral argument. The forcing in Theorem 5.5 is orthogonal to the Killing kernel, and division by positive eigenvalues gives the stated solution on its orthogonal complement.\par
The Section 6 Jacobi decomposition was checked with its spectral hypotheses. For \(v=x^{\top}Ax-2\operatorname{Tr}A\), where \(A\) is real symmetric, normalized Gaussian expectation with covariance \(2I\) gives \[\mathbb E[v^2]=8\operatorname{Tr}(A^2),\qquad
 q=|\nabla v|^2-\Delta|\nabla v|^2,\qquad
 \mathbb E[q^2]=128\operatorname{Tr}(A^4).\] Thus the quadratic obstruction cannot disappear through a trace-free or indefinite choice of \(A\). An independent warped-product calculation for \(g_t=(1+tu)g_N+g_{\mathbb R^{n-\ell}}\), \(f_t=\bar f+t\ell u/2\), gives \[\phi''_N(0)=-|\nabla u|^2g_N,\qquad
 \phi''_E(0)=-\ell u\,\operatorname{Hess}u-\tfrac\ell2\,du\otimes du.\] This verifies (8.2) using only \(\operatorname{Ric}_N=g_N/2\), without assuming constant sectional curvature. Symbolic residuals and twelve exact Gaussian matrix cases were saved in the native calculation record; the general argument, rather than those finite tests alone, supports the conclusion.\par
The gauge iteration is bounded in \(C^{5,\alpha}\) and contracting in \(C^{4,\alpha}\). Compactness and lower semicontinuity preserve the stronger bound in the limit. The center-of-mass adjustment treats the translation kernel, and the Gaussian cutoff-shell errors provide the improvement used above. The weighted \(L^1\) residual term must be retained from Theorem 9.1 and the preceding derivation; its omission in the last display (9.36) is not used to infer an additional defect.\par
Several local signs, radius labels and norm exponents are uniquely corrected by their immediately adjacent identities. They are not counted as independent Problems. Two claims remain bounded Needs Evidence items: the exact forcing coefficient in (5.9), which is stronger than the direct application of Theorem 4.2 records, and the unrestricted geometric scope of Proposition 5.8’s uniform elliptic estimates. The cylinder instance needed in Section 10 has the required geometric control. No counterexample to either broader statement was established, and these open items do not increase the report’s assessment.\par
\section*{Overall impact and state}
The highest supported assessment is Major because the missing relative endpoint hypothesis affects the core almost-splitting derivation. The boundary-term correction and parameter choices above support a weaker propagation estimate and the qualitative iteration. They do not certify the original fixed rate, a complete revised manuscript, or author adoption. The separate false integrated estimate has a fully verified local correction and is Minor. No central rigidity counterexample has been established.\par
Assessment | Repair verification | Adoption | Remaining work\par
Problem 1: Major | Boundary estimate proved; reduced-domain positive rate and subsequent qualitative iteration checked. | Unknown | Original fixed propagation rate not restored.\par
Problem 2: Minor | Correct coefficient proved by integration; counterexample verified. | Unknown | No independent acceptance.\par
Scope items: Needs Evidence | Cylinder application distinguished from general claims. | Unknown | Exact (5.9) coefficient and broad (5.8) scope unresolved.\par
Scientific decision and adjudication are pending. The human-review list is empty. Initial delivery is distinct from independent scientific acceptance.\par
9 T. H. Colding and W. P. Minicozzi II, Singularities of Ricci flow and diffeomorphisms, Publ. Math. IHÉS 142 (2025), 75–152. doi:10.1007/s10240-025-00157-1 (https://doi.org/10.1007/s10240-025-00157-1). Reviewed publisher version of record. C. Mantegazza and R. Müller, Perelman’s entropy functional at Type I singularities of the Ricci flow, arXiv:1205.4143v3 (https://arxiv.org/abs/1205.4143v3), 30 November 2012. Theorem 1.4 and adjoining definitions and remarks consulted; equivalence to the later publisher version is not asserted.\par
\section*{Supporting checks}
The canonical PDF SHA256 is \[\texttt{ec59dcce04edfd001c05dbda06c8cbd1}\] \[\texttt{508f88afb6e121adaf518ea12e6a8e78}.\] The matching verification directory retains the full source manifest, hash checks, text-reading intervals, actual source-image observations, analytical checkpoints, exact calculation inputs and outputs, current TeX/PDF correspondence, real compiler and dependency logs, and actual checks of every final report page. No author program or author-source compilation was used.\par
The root role integrated the scientific review and wrote this report. These are actual AI cross-checks within the current conversation, not human reviews or independent acceptance. No independent UI observation or backend equivalence is asserted.\par
\end{document}
